\documentclass[11pt,a4paper]{article}
\usepackage[margin=27mm]{geometry}
\usepackage[T1]{fontenc}
\usepackage{lmodern}
\usepackage{microtype}
\usepackage{amsmath,amssymb,amsthm,mathtools}
\usepackage{graphicx}
\usepackage[hidelinks]{hyperref}

\newtheorem{theorem}{Theorem}
\newtheorem{proposition}{Proposition}
\newcommand{\ellp}{\ell^p}
\newcommand{\ellq}{\ell^q}
\newcommand{\E}{\mathcal E}
\newcommand{\D}{\mathcal D}

\title{Recurrence does not control $\partial^2$ on $\ell^p$ when $p>2$}
\author{Autonomous proof-discovery packet for arXiv:2605.21069}
\date{11 August 2026}

\begin{document}
\maketitle

\noindent\textbf{Status.} Candidate full counterexample and sharpened positive replacement; likely valid, human review requested. Model: GPT5.6.

\section{Source question}

Bartmann and Keller study the boundary operator on weighted, possibly infinite simplicial complexes and characterize $\partial\partial=0$ on $\ell^2$ by recurrence of link graphs \cite{BK}. In the remark after their Theorem~4.1 (source PDF page~9), they prove the same implication on $\ell^p$ for $1\leq p\leq2$, observe that their argument fails for $p>2$, and ask for conditions governing the larger classes. The decisive excerpt is reproduced below.

\begin{center}
\includegraphics[width=\textwidth]{figures/open_problem_crop.png}
\end{center}

The result below gives a negative answer to recurrence sufficiency for every $p>2$. It also identifies the direct potential-theoretic replacement suggested by H\"older duality.

\section{Definitions and proof intuition}

For a fixed simplex $\rho$, write $b_\rho$ for its weighted link graph. If $1<q<\infty$, call a connected component $X$ of this link \emph{$q$-parabolic} when there are $0\leq\varphi_n\leq1$ in $C_c(X)$ such that $\varphi_n\to1$ pointwise and
\[
  \E_{q,\rho}(\varphi_n)
  :=\frac12\sum_{v,w\in X}b_\rho(v,w)
       |\varphi_n(v)-\varphi_n(w)|^q\longrightarrow0.
\]
This is the standard null-sequence formulation of graph $q$-parabolicity; see \cite{AFS} for the general locally summable setting.

\subsection*{Proof intuition}

The source's recurrence proof tests the total boundary against cutoffs tending to one. For an $\ell^p$ cochain, H\"older's inequality pairs it with the $\ell^q$ norm of the cutoff gradient, where $q=p/(p-1)$. Thus the exact hypothesis is not small quadratic energy unless $p=2$, but small $q$-energy. When $p>2$, one has $q<2$, and recurrence need not imply $q$-parabolicity.

The counterexample makes this mismatch visible as a conserved flux. Put conductance $n$ on the $n$th edge of a half-line and value $1/n$ on that edge. The weighted flux is identically one. Its divergence vanishes at every interior vertex but leaves one unit of charge at the endpoint. Quadratic resistance is $\sum 1/n=\infty$, so the link is recurrent; meanwhile the cochain's $\ell^p$ cost is $\sum n^{1-p}<\infty$ exactly when $p>2$.

\section{The positive replacement}

\begin{proposition}[Dual parabolicity suffices]\label{prop:positive}
Let $1<p<\infty$ and $q=p/(p-1)$. Fix $\rho$ in a weighted simplicial complex in the sense of \cite{BK}. If every connected component of $b_\rho$ is $q$-parabolic, then
\[
  \partial\partial\omega(\rho)=0
  \qquad
  \bigl(\omega\in D(\partial\partial)\cap\ell^p(\Sigma,m)\bigr).
\]
\end{proposition}

\begin{proof}
Enumerate the components of $b_\rho$. By their $q$-parabolic null sequences and a diagonal choice, there are $0\leq\varphi_n\leq1$ in $C_c(\operatorname{lk}(\rho))$ such that $\varphi_n\to1$ pointwise and $\E_{q,\rho}(\varphi_n)\to0$. Indeed, on the first $n$ components choose cutoffs with total $q$-energy at most $1/n$, and take their sum; monotone componentwise null sequences give pointwise convergence.

Use the localization maps $\pi_\rho,\pi^\rho$ of \cite{BK}. Since $\partial\omega\in D(\partial)$, the function $\pi_\rho\partial\omega$ is in $\ell^1(\operatorname{lk}(\rho),m_\rho)$. Dominated convergence and the finite-support Stokes identity therefore give
\begin{align*}
 m(\rho)\,|\partial\partial\omega(\rho)|
 &=\lim_{n\to\infty}
   \left|\sum_{\operatorname{lk}(\rho)}
       m_\rho(\pi_\rho\partial\omega)\varphi_n\right|\\
 &=\lim_{n\to\infty}
   \left|\sum_{\Sigma}m\omega\,\delta\pi^\rho\varphi_n\right|\\
 &\leq \|\omega\|_p
       \lim_{n\to\infty}\|\delta\pi^\rho\varphi_n\|_q.
\end{align*}
For every simplex $vv'\rho$ containing two link vertices, the only two terms in the localized coboundary differ by orientation and give $\varphi_n(v)-\varphi_n(v')$. Hence, exactly
\[
 \|\delta\pi^\rho\varphi_n\|_q^q
 =\frac12\sum_{v,v'}b_\rho(v,v')
       |\varphi_n(v)-\varphi_n(v')|^q
 =\E_{q,\rho}(\varphi_n)\longrightarrow0.
\]
The displayed boundary of boundary is therefore zero.
\end{proof}

\section{A recurrent counterexample for every \texorpdfstring{$p>2$}{p greater than 2}}

\begin{theorem}\label{thm:counterexample}
There is one weighted one-dimensional simplicial complex $(\Sigma,m)$ such that every link component is recurrent and, simultaneously for every $p>2$, there is a single cochain
\[
 \omega\in D(\partial\partial)\cap\ell^p(\Sigma,m)
 \quad\text{with}\quad
 \partial\partial\omega(\varnothing)=-1.
\]
Consequently, recurrence of the links does not imply the complex property on $\ell^p$ for any $p>2$.
\end{theorem}

\begin{proof}
Let the vertices be $v_0,v_1,\ldots$ and let
$e_n=\{v_{n-1},v_n\}$ for $n\geq1$. Take the augmented one-dimensional complex
\[
 \Sigma=\{\varnothing\}\cup\{v_n:n\geq0\}
       \cup\{e_n:n\geq1\}
\]
with weights
\[
 m(\varnothing)=1,
 \qquad m(v_n)=2^{-n-1},
 \qquad m(e_n)=n.
\]
This is locally summable in the source's sense: $\sum_n m(v_n)=1$ at the empty simplex, and each vertex has only one or two edge cofaces of finite weight.

Orient the augmented coboundary by
\[
 \theta(\varnothing,v_n)=1,
 \qquad \theta(v_{n-1},e_n)=-1,
 \qquad \theta(v_n,e_n)=1.
\]
Then $\delta^2=0$, since the two vertex contributions on each edge cancel. Define $\omega$ to vanish on the empty simplex and all vertices and set
\[
 \omega(e_n)=\frac1n\qquad(n\geq1).
\]
For every $p>2$,
\[
 \|\omega\|_p^p
 =\sum_{n=1}^\infty m(e_n)|\omega(e_n)|^p
 =\sum_{n=1}^\infty n^{1-p}<\infty.
\]

All sums defining $\partial\omega$ at vertices are finite. At the endpoint and the interior vertices, respectively,
\[
 \partial\omega(v_0)
 =\frac{-m(e_1)\omega(e_1)}{m(v_0)}=-2,
 \qquad
 \partial\omega(v_n)
 =\frac{m(e_n)\omega(e_n)-m(e_{n+1})\omega(e_{n+1})}{m(v_n)}=0
 \quad(n\geq1).
\]
It vanishes on edges. Thus $\partial\omega\in D(\partial)$: the only nontrivial formal-domain sum is
\[
 \sum_{n\geq0}m(v_n)|\partial\omega(v_n)|=1.
\]
Therefore $\omega\in D(\partial\partial)$, and
\[
 \partial\partial\omega(\varnothing)
 =\frac1{m(\varnothing)}
   \sum_{n\geq0}m(v_n)\theta(\varnothing,v_n)
          \partial\omega(v_n)=-1.
\]

It remains to check the link hypothesis. The link of $\varnothing$ is the half-line with conductance
$b_\varnothing(v_{n-1},v_n)=m(e_n)=n$. Let $H_N=\sum_{j=1}^N1/j$ and define a finitely supported cutoff by
\[
 \eta_N(v_k)=
 \begin{cases}
  1-H_k/H_N,&0\leq k<N,\\
  0,&k\geq N,
 \end{cases}
 \qquad H_0=0.
\]
Then $\eta_N\to1$ pointwise and
\[
 \E_{2,\varnothing}(\eta_N)
 =\sum_{n=1}^N n\left(\frac1{nH_N}\right)^2
 =\frac1{H_N}\longrightarrow0.
\]
Thus this link is $2$-parabolic, equivalently recurrent. Links of vertices have no link edges, so all their components are trivially recurrent; higher links are empty.

For completeness, this example also pinpoints Proposition~\ref{prop:positive}. If $q=p/(p-1)<2$ and $\phi\in C_c$ satisfies $\phi(v_0)=1$, telescoping and weighted H\"older give
\[
 1
 \leq
 \left(\sum_{n\geq1}n|\phi(v_{n-1})-\phi(v_n)|^q\right)^{1/q}
 \left(\sum_{n\geq1}n^{-(p-1)}\right)^{1/p}.
\]
The $q$-energy of such cutoffs is bounded away from zero. Hence the recurrent half-line is $q$-hyperbolic for precisely the dual exponents relevant to every $p>2$.
\end{proof}

\section{Scope, novelty, and review recommendation}

Theorem~\ref{thm:counterexample} answers the source's $p>2$ direction negatively if recurrence is retained as the proposed condition. Proposition~\ref{prop:positive} supplies a general sufficient condition but does not claim an unrestricted converse for arbitrary higher-dimensional links; extra faces can affect the global domain of a cochain built from a link flow. No $\ell^p$-Hodge decomposition or Riesz-transform equivalence is asserted.

The bounded novelty check (through 11 August 2026) searched all four local run indexes, the exact arXiv id and title, and combinations of \emph{boundary of boundary}, \emph{complex property}, $\ell^p$, $p>2$, \emph{weighted half-line}, \emph{unit flux}, and $q$-parabolicity. It found the source and the graph-parabolicity background \cite{AFS}, but no paper stating this counterexample or applying dual parabolicity to the source question. Because arXiv:2605.21069 is very recent, novelty is provisional.

Human review should focus on the iterated formal-domain check, the augmented orientation at $\varnothing$, the componentwise diagonal cutoff in Proposition~\ref{prop:positive}, and a refreshed literature search.

\begin{thebibliography}{9}
\bibitem{BK}
P.~Bartmann and M.~Keller,
\emph{The complex property of the boundary operator on simplicial complexes},
arXiv:2605.21069v1, 2026.

\bibitem{AFS}
A.~Adriani, F.~Fischer, and A.~G.~Setti,
\emph{Characterizations of $p$-Parabolicity on Graphs},
arXiv:2507.13696v2, 2026.
\end{thebibliography}

\end{document}
